% Copyright 2026 Myeongjin Kang. Original prose: CC BY 4.0
% (https://creativecommons.org/licenses/by/4.0/). Statements quoted from
% third-party works remain under their authors' rights.
\documentclass[11pt,a4paper]{article}
\usepackage[T1]{fontenc}
\usepackage{lmodern}
\usepackage[margin=1in]{geometry}
\usepackage{amsmath,amssymb,amsthm,mathtools}
\usepackage{microtype}
\usepackage{booktabs,array}
\usepackage{enumitem}
\usepackage{xcolor}
\usepackage[colorlinks=true,linkcolor=blue!50!black,citecolor=blue!50!black,urlcolor=blue!50!black]{hyperref}
\hypersetup{pdftitle={The crossing number of the complete graph minus an edge},
  pdfauthor={Myeongjin Kang},
  pdfsubject={Crossing numbers; complete graphs; Cartesian products with paths},
  pdfkeywords={crossing number, complete graph, edge deletion, Zheng conjecture, Cartesian product, path}}
\setlength{\emergencystretch}{2em}
\newcolumntype{P}[1]{>{\raggedright\arraybackslash}p{#1}}
\newcommand{\doi}[1]{\href{https://doi.org/#1}{doi:#1}}

\newtheorem{maintheorem}{Theorem}
\renewcommand{\themaintheorem}{\Alph{maintheorem}}
\newtheorem{maincorollary}[maintheorem]{Corollary}
\newtheorem{premise}{Premise}
\renewcommand{\thepremise}{U\arabic{premise}}
\newtheorem{theorem}{Theorem}[section]
\newtheorem{lemma}[theorem]{Lemma}
\newtheorem{proposition}[theorem]{Proposition}
\newtheorem{corollary}[theorem]{Corollary}
\theoremstyle{definition}
\newtheorem{definition}[theorem]{Definition}
\theoremstyle{remark}
\newtheorem{remark}[theorem]{Remark}
\numberwithin{equation}{section}

\newcommand{\R}{\mathbb R}
\newcommand{\crn}{\operatorname{cr}}
\newcommand{\PP}{\mathcal P}
\newcommand{\SSS}{\mathcal S}
\newcommand{\TT}{\mathcal T}
\newcommand{\VV}{\mathcal V}
\newcommand{\AAA}{\mathcal A}
\newcommand{\ee}{\mathbf e}
\newcommand{\Sym}{\mathrm{Sym}}
\newcommand{\tg}{\tau}

\title{The crossing number of the complete graph minus an edge}
\author{Myeongjin Kang}
\date{October 2026}

\begin{document}
\maketitle

\begin{abstract}
We prove that for every integer $n\ge 2$ and every edge $e$ of the complete graph $K_n$, the ordinary crossing number of $K_n-e$ is
\[
\crn(K_n-e)=\frac14\Big\lfloor\frac{n+2}{2}\Big\rfloor\Big\lfloor\frac{n-1}{2}\Big\rfloor\Big\lfloor\frac{n-3}{2}\Big\rfloor\Big\lfloor\frac{n-4}{2}\Big\rfloor ,
\]
which was conjectured by Zheng, Lin, Yang and Cui in 2007. The proof relies on the September~2026 OpenAI preprint proving the Harary--Hill formula for $\crn(K_n)$, whose stated results we use without re-proving them. Given that formula, the even orders follow from a theorem of Ouyang, Wang and Huang (2014), equivalently of Chia and Lee (2015), and the odd orders $n\le 13$ follow from their counting bound and integrality. The new content is the lower bound for odd $n\ge 15$. We adapt the polynomial rank method of the OpenAI preprint to the missing edge. The adaptation restricts the test polynomials so that they put no weight on the missing pair. It pairs two different vanishing subspaces by mixed orthogonality. It also localizes, with a twisted detector that needs the full rectangular nonsingularity statement of the preprint, only at the vertices whose stars are complete. Via the product identity of Ouyang, Wang and Huang, the formula gives the crossing number of $K_m\,\square\,P_n$ conjectured by Zheng et al.\ for all $m\ge 4$ and $n\ge 1$, where $P_n$ is the path with $n$ edges. This consequence is new for odd $m\ge 13$.
\end{abstract}

\section{Introduction}\label{sec:intro}

An \emph{ordinary drawing} of a finite simple graph places its vertices at distinct points of the plane and its edges as simple arcs joining their endpoints. Edge interiors avoid the vertices, and two edges meet in finitely many proper crossings, with no tangencies or triple points. The number $c(D)$ of crossings of a drawing $D$ counts every crossing point, so two edges that meet several times contribute several crossings. The \emph{crossing number} $\crn(G)$ is the minimum of $c(D)$ over all ordinary drawings of $G$. No restriction is placed on vertex positions or edge routes. This is the convention of~\cite[Section~1]{OAIKn}, and all statements below use it.

Write
\[
Z(n)=\frac14\Big\lfloor\frac n2\Big\rfloor\Big\lfloor\frac{n-1}2\Big\rfloor\Big\lfloor\frac{n-2}2\Big\rfloor\Big\lfloor\frac{n-3}2\Big\rfloor .
\]
The OpenAI preprint~\cite[Theorem~1.1]{OAIKn} proves that $\crn(K_n)=Z(n)$ for every $n\ge 3$, which is the Harary--Hill conjecture~\cite{HararyHill1963}. Our main result is the corresponding formula for the complete graph with one edge removed.

\begin{maintheorem}\label{thm:A}
For every integer $n\ge 2$ and every edge $e$ of $K_n$,
\begin{equation}\label{eq:mainformula}
\crn(K_n-e)=Z(n)-\binom{\lfloor (n-1)/2\rfloor}{2}
=\frac14\Big\lfloor\frac{n+2}{2}\Big\rfloor\Big\lfloor\frac{n-1}{2}\Big\rfloor\Big\lfloor\frac{n-3}{2}\Big\rfloor\Big\lfloor\frac{n-4}{2}\Big\rfloor .
\end{equation}
Equivalently, for every integer $s\ge 2$,
\begin{equation}\label{eq:parity}
\crn(K_{2s+1}-e)=\binom{s}{2}^{2}-\binom{s}{2},
\qquad
\crn(K_{2s}-e)=\binom{s}{2}\binom{s-1}{2}-\binom{s-1}{2}=\frac{(s^2-1)(s-2)^2}{4}.
\end{equation}
\end{maintheorem}

The two expressions in~\eqref{eq:mainformula} agree for every $n\ge2$. For $n=2s+1$ the four floors are $s+1,s,s-1,s-2$, and $\binom s2^2-\binom s2=\frac{s(s-1)(s+1)(s-2)}{4}$. For $n=2s$ they are $s+1,s-1,s-2,s-2$, and $Z(2s)-\binom{s-1}2=\frac{(s-1)(s-2)}{4}\bigl(s(s-1)-2\bigr)=\frac{(s+1)(s-1)(s-2)^2}{4}$. All automorphisms of $K_n$ act transitively on edges, so all graphs $K_n-e$ are isomorphic. The value is $0$ for $n\le5$.

For a graph $G$ and the path $P_n$ with vertices $0,1,\dots,n$ (so $n$ edges, the convention of~\cite[Section~5]{OWH}), $G\,\square\,P_n$ denotes the Cartesian product.

\begin{maincorollary}\label{cor:B}
For all integers $m\ge4$ and $n\ge1$,
\begin{equation}\label{eq:zheng}
\crn(K_m\,\square\,P_n)=\frac14\Big\lfloor\frac{m+1}{2}\Big\rfloor\Big\lfloor\frac{m-1}{2}\Big\rfloor\Big\lfloor\frac{m-2}{2}\Big\rfloor\Big(n\Big\lfloor\frac{m+4}{2}\Big\rfloor+\Big\lfloor\frac{m-4}{2}\Big\rfloor\Big).
\end{equation}
\end{maincorollary}

Corollary~\ref{cor:B} is one application of Theorem~\ref{thm:A}, through the identity $\crn(K_m\,\square\,P_n)=(n-1)\crn(K_{m+2}-e)+2\crn(K_{m+1})$ of Ouyang, Wang and Huang~\cite[Corollary~5.1]{OWH}. We do not regard it as a separate result.

\subsection{History and what is new}

Zheng, Lin, Yang and Cui~\cite{ZLYC} gave the upper bound~\eqref{eq:mainformula} for $\crn(K_n-e)$, conjectured equality, and proved it for $n\le 8$. In the same paper they proved $\crn(K_6\,\square\,P_n)=15n+3$, proved the lower bound $\crn(K_m\,\square\,P_n)\ge(n-1)\crn(K_{m+2}-e)+2\crn(K_{m+1})$, and conjectured~\eqref{eq:zheng}. We were not able to obtain the text of~\cite{ZLYC}. These attributions follow Ouyang, Wang and Huang~\cite[Section~1, display~(1.1)--(1.2), and Lemma~2.3]{OWH} and the survey of Clancy, Haythorpe and Newcombe~\cite[Theorem~2.19, Conjecture~2.20, Theorems~3.12--3.13, Conjecture~3.14]{CHN}.

Ouyang, Wang and Huang~\cite{OWH} proved two lower bounds for $\crn(K_m-e)$ in terms of complete-graph crossing numbers~\cite[Theorems~3.1 and~3.2]{OWH}. They deduced the even case of the conjecture from the Harary--Hill formula for the preceding odd order~\cite[Theorem~3.3]{OWH}, and settled $m\le 12$ using the then-known values of $\crn(K_m)$~\cite[Theorem~3.4]{OWH}. Chia and Lee~\cite{ChiaLee} also obtained equivalent results: the upper bound~\cite[Lemma~1]{ChiaLee}, the counting lower bound~\cite[Theorem~1]{ChiaLee}, the even case under the odd Harary--Hill hypothesis~\cite[Theorem~2]{ChiaLee}, the values for $n\le12$~\cite[Corollary~1]{ChiaLee}, and the conjecture~\cite[Conjecture~1]{ChiaLee}. Ouyang, Wang and Huang also proved the product identity~\cite[Corollary~5.1]{OWH} and deduced~\eqref{eq:zheng} for $m\le 10$~\cite[Corollary~5.2]{OWH}.

Once the Harary--Hill formula is available for every order~\cite{OAIKn}, the following parts of Theorem~\ref{thm:A} are \emph{classical consequences} and are not claimed as new:
\begin{itemize}[leftmargin=1.5em]
\item the upper bound for every $n$~\cite{ZLYC,ChiaLee}; Section~\ref{sec:upper} recounts it in the two-page drawing of~\cite{OAIKn};
\item the lower bound for every even $n$, by~\cite[Theorem~3.3]{OWH} or~\cite[Theorem~2]{ChiaLee}; Section~\ref{sec:even} reproduces the proof of~\cite[Theorem~3.2]{OWH};
\item the lower bound for odd $n\le 13$, from~\cite[Theorem~3.1]{OWH} (equivalently~\cite[Theorem~1]{ChiaLee}) and integrality; see Section~\ref{sec:novelty}.
\end{itemize}
The \emph{new} part is the lower bound $\crn(K_{2s+1}-e)\ge\binom s2^2-\binom s2$ for $s\ge7$, that is, for odd $n\ge 15$ (Proposition~\ref{prop:odd}). There the counting bound of~\cite[Theorem~3.1]{OWH} falls short of the target by exactly $\frac{(s-1)(s-2)}{2(2s+1)}$, which is at least $1$ for every $s\ge 7$. Our proof of Proposition~\ref{prop:odd} covers every odd $n\ge5$, but for $n\le13$ it only reproves known consequences. Correspondingly, Corollary~\ref{cor:B} is new exactly for odd $m\ge 13$. For even $m$ and odd $m\le11$, the value $\crn(K_{m+2}-e)$ needed in~\cite[Corollary~5.1]{OWH} was already available after the Harary--Hill formula.

\subsection{Method and relation to other work}

The OpenAI preprint~\cite{OAIKn} proves the odd lower bound $\crn(K_{2s+1})\ge\binom s2^2$ as follows. A skew bilinear form built from signed crossings and rotation orders vanishes on pairs of cycles. Weighting edges by values of interpolation-weighted polynomials yields cycles, hence polynomial identities. These identities make a space of polynomials of dimension $\binom s2^2$ embed isotropically into the coordinates indexed by crossing edge pairs.

For $K_n-e$ the same polynomial chains put nonzero weight on the missing pair. Four changes are needed, and they are the content of Section~\ref{sec:odd}.
\begin{enumerate}[leftmargin=1.7em,label=(\roman*)]
\item \emph{Missing-pair tests} (Lemma~\ref{lem:tests}). A test polynomial gives a valid identity whenever it satisfies two polynomial symmetry conditions at the missing pair. The conditions are sufficient, not necessary (Remark~\ref{rem:nonconverse}).
\item \emph{Split spaces and mixed orthogonality} (Lemmas~\ref{lem:dim} and~\ref{lem:mixed}). The full space of~\cite{OAIKn} is not admissible. We use two subspaces $\SSS_a,\SSS_b$, each of codimension $\binom s2$, vanishing on the missing pair in the first or in the second edge slot, and show that their images are orthogonal to each other. Their intersection would lose $\binom s2-1$ further dimensions per side.
\item \emph{Twisted detector} (Lemma~\ref{lem:twisted}). The localization tests must carry extra factors $(y-t_a)$ and $(v-t_a)$. The resulting pairing is controlled by the \emph{rectangular} nonsingularity in~\cite[Lemma~3.3]{OAIKn}, not by its square consequence.
\item \emph{Localization away from the missing edge} (Lemma~\ref{lem:injective}). Only the $2s-1$ vertices with complete stars are used for localization. This still exceeds the degree $2s-4$ of the relevant diagonal polynomial, and the degree descent preserves the vanishing condition at the missing pair.
\end{enumerate}

Sadhu~\cite{Sadhu} recently determined $\crn(K_{m,n}-e)$ for complete bipartite graphs, a family also treated by Chia and Lee~\cite{ChiaLee}. That argument is stated to be inspired by the same algebraic framework, in its bipartite form~\cite{OAIKmn}. That preprint appeared before the present work. Its stated new ingredients are exact node-deletion maps, a balanced star-compression rank inequality with dimension losses, an extension of the bipartite order detector to unequal rectangles, and a signed three-edge completion that replaces the missing edge. In the passages we compared (Section~\ref{sec:novelty}), these operations differ from ours. Neither theorem implies the other. We therefore claim the specific argument for $K_n-e$, not the general idea of adapting the framework of~\cite{OAIKn} to a deleted edge.

Asiri, Burdett, Chimani, Haythorpe, Newcombe and Wagner~\cite{Asiri} gave a general computational framework for $\crn(G\,\square\,P_n)$ and applied it to graphs of order five and six. It does not treat complete graphs of arbitrary order.

\subsection{Organization and status}
Section~\ref{sec:premises} lists the exact upstream statements used. Section~\ref{sec:odd} proves the odd lower bound. Section~\ref{sec:even} reproduces the even lower bound. Section~\ref{sec:upper} gives the upper bound and Section~\ref{sec:proofs} assembles Theorem~\ref{thm:A} and Corollary~\ref{cor:B}. Section~\ref{sec:novelty} records the novelty boundary, Section~\ref{sec:computations} the computational support, and Section~\ref{sec:formal} the status of formal verification, which is pending. The finite computations support the proof but are not part of it.

\section{Conventions and upstream statements}\label{sec:premises}

We use the following statements from the OpenAI preprint~\cite{OAIKn}, as fixed in its source at commit \texttt{fd4aeeb2ee4fc729c18d98444fed42fd0529eeeb} of \url{https://github.com/openai/math}, file \path{preprints/The-crossing-number-of-complete-graphs-September-23-2026/build/main.tex} (SHA-256 \texttt{5ae2c607\ldots2301}). We cite them by their source labels and use them as stated. We do not re-prove them, and the correctness of our results depends on theirs.

Let $G$ be a finite simple graph with a fixed orientation of each edge. Its real edge space $C_1(G)$ has the oriented edges as a basis; reversing an orientation negates the vector. For a vertex set $\{1,\dots,n\}$ we write $[ip]$ for the oriented pair from $i$ to $p$, with $[pi]=-[ip]$ and $[ii]=0$. The boundary map is $\partial[ip]=\ee_p-\ee_i$, and a \emph{cycle} is an element of $\ker\partial$. Following~\cite[Section~2]{OAIKn}, a drawing is \emph{normalized} if it is polygonal, straight near every vertex and crossing, and has no crossing between two edges sharing an endpoint. Independent edges may cross repeatedly.

\begin{premise}[{\cite[Lemma~A.1, label \texttt{lem:normalization}]{OAIKn}}]\label{U:norm}
For every finite simple graph, a minimum ordinary drawing can be chosen normalized.
\end{premise}

For a normalized drawing, let $I(e,f)$ be the sum of the signs of the interior crossings of two distinct oriented edges, with $I(e,e)=0$. At each vertex $w$, cut the circular order of the outgoing spokes at a gap and list the incident edges counterclockwise. Put $T_w(e,f)=1$ if $e$ precedes $f$, $-1$ if $f$ precedes $e$, and $0$ if $e=f$. Put $\sigma_{we}=1$ if $e$ points out of $w$ and $-1$ if it points into $w$. Define a bilinear form on $C_1(G)$ by
\begin{equation}\label{eq:J}
J(e,f)=I(e,f)+\frac12\sum_{w\text{ a common endpoint of }e,f}\sigma_{we}\sigma_{wf}T_w(e,f)
\end{equation}
\cite[equation \texttt{eq:J}]{OAIKn}. It is skew-symmetric.

\begin{premise}[{\cite[Lemma~2.1, label \texttt{lem:cycles}]{OAIKn}}]\label{U:cycles}
For a normalized drawing of any finite simple graph $G$, $J(z,z')=0$ for any two cycles $z,z'$ of $G$.
\end{premise}

The proof in~\cite{OAIKn} states explicitly that the argument applies to every finite simple graph; we apply it to $G=K_n-e$, never to $K_n$.

\begin{premise}[{\cite[Lemma~3.1, label \texttt{lem:weights}]{OAIKn}}]\label{U:weights}
For distinct reals $u_1,\dots,u_N$ and $\Gamma(t)=\prod_{p=1}^N(t-u_p)$, one has $\sum_{p=1}^N A(u_p)/\Gamma'(u_p)=0$ for every polynomial $A$ of degree at most $N-2$.
\end{premise}

For nonnegative integers $\alpha,\beta$ let $\PP_{\alpha,\beta}$ be the real polynomials in $(u,v)$ of degree at most $\alpha$ in $u$ and at most $\beta$ in $v$. Let $s\ge2$ and let $u_1,\dots,u_{2s}$ be distinct reals carrying an arbitrary prescribed linear order, not necessarily their numerical order. Put $S(p,q)=\frac12$ if $p$ precedes $q$, $-\frac12$ if $q$ precedes $p$, $S(p,p)=0$, $\Gamma(t)=\prod_{p=1}^{2s}(t-u_p)$, and
\begin{equation}\label{eq:M}
M(A,B)=\sum_{p,q=1}^{2s}S(p,q)\,\frac{u_p-u_q}{\Gamma'(u_p)\Gamma'(u_q)}\,A(u_p,u_q)\,B(u_p,u_q)
\end{equation}
\cite[equation \texttt{eq:M}]{OAIKn}.

\begin{premise}[{\cite[Lemma~3.3, label \texttt{lem:detector}, first sentence]{OAIKn}}]\label{U:detector}
The pairing $M$ on $\PP_{s-1,s-2}\times\PP_{s-2,s-1}$ is nonsingular.
\end{premise}

\begin{premise}[{\cite[Theorem~1.1, label \texttt{thm:main}]{OAIKn}}]\label{U:HH}
$\crn(K_n)=Z(n)$ for every integer $n\ge3$.
\end{premise}

\begin{premise}[{\cite[Section~5 and Proposition~5.1, labels \texttt{sec:upper}, \texttt{prop:upper}]{OAIKn}}]\label{U:upper}
Let $n\ge3$ and $m=\lfloor n/2\rfloor$. Place vertices $0,1,\dots,n-1$ at increasing positions on a line, and draw the edge $ij$ as a semicircle above the line if $(i+j)\bmod n<m$ and below it otherwise. For generic positions this is an ordinary drawing of $K_n$. Its crossings are exactly the pairs of edges on the same side whose four endpoints alternate along the line, each such pair crossing once, and it has exactly $Z(n)$ crossings.
\end{premise}

The bipartite companion~\cite{OAIKmn} is not used. In~\cite{OAIKn}, the results labelled \texttt{lem:tests}, \texttt{lem:injectivity} and \texttt{prop:odd} concern $K_n$. We do not use them as black boxes. Their proofs serve as templates for Lemmas~\ref{lem:tests}, \ref{lem:injective} and Proposition~\ref{prop:odd}, and every changed step is proved below.

The external results used are as follows.
\begin{itemize}[leftmargin=1.5em]
\item \cite[Corollary~5.1]{OWH}: \emph{for $m\ge4$ and $n\ge1$, $\crn(K_m\,\square\,P_n)=(n-1)\crn(K_{m+2}\setminus e)+2\crn(K_{m+1})$}, with $P_n$ the path with vertices $0,\dots,n$. It is unconditional. Its proof uses~\cite[Theorems~5.1, 5.2 and Lemma~4.1]{OWH} and a lemma of Bokal~\cite{Bokal2007}. We cite it and have not re-checked that chain.
\item \cite[Theorem~3.2]{OWH} is reproduced with proof in Section~\ref{sec:even}. \cite[Theorem~3.1]{OWH} is used only in Section~\ref{sec:novelty} to delimit the novelty.
\end{itemize}

Table~\ref{tab:deps} summarizes where each input enters.

\begin{table}[ht]
\centering
\small
\begin{tabular}{@{}P{0.25\textwidth} P{0.33\textwidth} P{0.36\textwidth}@{}}
\toprule
Input & Source locator & Used in \\
\midrule
U1 normalization & \cite{OAIKn}, \texttt{lem:normalization} & Odd (Sec.~\ref{sec:odd}) and even (Sec.~\ref{sec:even}) lower bounds \\
U2 cycle pairing & \cite{OAIKn}, \texttt{eq:J}, \texttt{lem:cycles} & Lemma~\ref{lem:tests}, applied to $G=K_n-e$ \\
U3 interpolation weights & \cite{OAIKn}, \texttt{lem:weights} & Lemma~\ref{lem:tests} (boundary of test chains) \\
U4 rectangular detector & \cite{OAIKn}, \texttt{eq:M}, \texttt{lem:detector} & Lemma~\ref{lem:twisted} \\
U5 Harary--Hill formula & \cite{OAIKn}, \texttt{thm:main} & Even lower bound; Corollary~\ref{cor:B} \\
U6 two-page drawing & \cite{OAIKn}, \texttt{sec:upper}, \texttt{prop:upper} & Upper bound (Sec.~\ref{sec:upper}) \\
OWH Cor.~5.1 & \cite{OWH}, Corollary 5.1 (cited) & Corollary~\ref{cor:B} \\
OWH Thm.~3.2 & \cite{OWH}, Theorem 3.2 (reproved) & Even lower bound (Sec.~\ref{sec:even}) \\
OWH Thm.~3.1 & \cite{OWH}, Theorem 3.1 & Novelty boundary only (Sec.~\ref{sec:novelty}) \\
\bottomrule
\end{tabular}
\caption{Dependencies. The truth of U1--U6 is taken from~\cite{OAIKn}; OWH Corollary~5.1 is cited from the published paper.}
\label{tab:deps}
\end{table}

\section{The odd lower bound}\label{sec:odd}

Throughout this section $s\ge3$ and $n=2s+1$. Let $G=K_n-ab$ for two distinct vertices $a,b$ of $K_n$, and fix a normalized minimum drawing $D$ of $G$ (Premise~\ref{U:norm}), so $c(D)=\crn(G)$. Fix an orientation of each edge of $G$ and let $J$ be the form~\eqref{eq:J} on $C_1(G)$.

\subsection{Zero extension and test chains}
Regard $C_1(G)$ as the subspace of $C_1(K_n)$ spanned by all oriented pairs except $\pm[ab]$. Extend $J$ to a bilinear form $\widetilde J$ on $C_1(K_n)$ by setting $\widetilde J(f,g)=0$ whenever $f$ or $g$ is $[ab]$. Then $\widetilde J$ is skew-symmetric and restricts to $J$ on $C_1(G)$. For ordered vertex pairs we write $J([ip],[kq])$ for $\widetilde J([ip],[kq])$; it vanishes when $i=p$, when $k=q$, or when $\{i,p\}$ or $\{k,q\}$ equals $\{a,b\}$.

Choose distinct reals $t_1,\dots,t_n$ and set
\[
F(t)=\prod_{i=1}^n(t-t_i),\qquad w_i=\frac1{F'(t_i)}\neq0.
\]
For a polynomial $H(x,u,y,v)$ write
\begin{equation}\label{eq:testsum}
\Sigma(H)=\sum_{i,p,k,q=1}^{n}J([ip],[kq])\,w_iw_pw_kw_q\,H(t_i,t_p,t_k,t_q).
\end{equation}

\begin{lemma}[Missing-pair cycle tests]\label{lem:tests}
Let $H(x,u,y,v)$ have degree at most $n-2$ in each variable separately and satisfy the polynomial identities
\begin{equation}\label{eq:testconditions}
H(t_a,t_b,y,v)=H(t_b,t_a,y,v),\qquad H(x,u,t_a,t_b)=H(x,u,t_b,t_a).
\end{equation}
Then $\Sigma(H)=0$.
\end{lemma}

\begin{proof}
\emph{Tensor decomposition.} Let $\VV$ be the space of polynomials $A(x,u)$ of degree at most $n-2$ in each of $x,u$. Multiplication $A\otimes B\mapsto A(x,u)B(y,v)$ identifies $\VV\otimes\VV$ with the space of four-variable polynomials of separate degree at most $n-2$, since it maps a monomial basis to a monomial basis. Let $\varphi(A)=A(t_a,t_b)-A(t_b,t_a)$. This linear functional is nonzero, since $\varphi(x)=t_a-t_b$. Choose $e_0\in\VV$ with $\varphi(e_0)=1$ and a basis $e_1,\dots,e_N$ of $\AAA=\ker\varphi$. Write $H=\sum_{j,k=0}^N c_{jk}\,e_j(x,u)e_k(y,v)$. The first identity in~\eqref{eq:testconditions} says $(\varphi\otimes\mathrm{id})H=\sum_k c_{0k}e_k(y,v)=0$, so $c_{0k}=0$ for all $k$. The second says $(\mathrm{id}\otimes\varphi)H=\sum_jc_{j0}e_j(x,u)=0$, so $c_{j0}=0$ for all $j$. Hence $H=\sum_r A_r(x,u)B_r(y,v)$ with all $A_r,B_r\in\AAA$.

\emph{Test chains are cycles of $G$.} For $A\in\VV$ put
\[
z_A=\sum_{i,p=1}^nw_iw_p\,A(t_i,t_p)\,[ip]\in C_1(K_n).
\]
Its coefficient on the basis vector $[ab]$ is $w_aw_b\bigl(A(t_a,t_b)-A(t_b,t_a)\bigr)=w_aw_b\varphi(A)$, because the two ordered terms carry opposite vectors $[ab]$ and $[ba]=-[ab]$. Thus $z_A\in C_1(G)$ whenever $A\in\AAA$. For a monomial $A=x^\alpha u^\beta$ with $\alpha,\beta\le n-2$,
\[
\partial z_A=\Bigl(\sum_iw_it_i^\alpha\Bigr)\sum_pw_pt_p^\beta\ee_p-\Bigl(\sum_pw_pt_p^\beta\Bigr)\sum_iw_it_i^\alpha\ee_i=0
\]
by Premise~\ref{U:weights} with $N=n$. By linearity $\partial z_A=0$ for all $A\in\VV$. The boundary map of $G$ is the restriction of that of $K_n$, so $z_A$ is a cycle of $G$ for every $A\in\AAA$.

\emph{Conclusion.} For $A,B\in\AAA$, bilinearity gives $\Sigma\bigl(A(x,u)B(y,v)\bigr)=\widetilde J(z_A,z_B)=J(z_A,z_B)$, which is $0$ by Premise~\ref{U:cycles} applied to $G$. Summing over $r$ gives $\Sigma(H)=0$.
\end{proof}

\begin{remark}[The conditions are sufficient only]\label{rem:nonconverse}
The identities~\eqref{eq:testconditions} are not necessary for $\Sigma(H)=0$, and we make no claim that they characterize the vanishing tests. For example, $H=xy$ violates both identities, since $H(t_a,t_b,y,v)-H(t_b,t_a,y,v)=(t_a-t_b)y$ and $H(x,u,t_a,t_b)-H(x,u,t_b,t_a)=(t_a-t_b)x$. Yet $\Sigma(xy)=\widetilde J(z_x,z_x)=0$ by skew-symmetry, for every drawing. The proof below uses only the sufficient direction.
\end{remark}

\subsection{Support coordinates and split spaces}
Let $\TT$ be the set of ordered pairs $(f,g)$ of (undirected) edges of $G$ with no common endpoint and $J(f,g)\ne0$. The condition does not depend on the chosen orientations, and $\TT$ is closed under $(f,g)\mapsto(g,f)$ because $J$ is skew. For independent edges the correction term in~\eqref{eq:J} is empty, so $J(f,g)=I(f,g)\ne0$ forces $f$ and $g$ to cross. Each crossing involves exactly two edges, so
\begin{equation}\label{eq:support}
c(D)\ge\frac{|\TT|}{2}.
\end{equation}
Repeated crossings only increase $c(D)$, and cancellation of signs only removes pairs from $\TT$.

Let $\SSS$ be the space of polynomials $P(x,u,y,v)$ of degree at most $s-2$ in each variable that are symmetric under $x\leftrightarrow u$ and under $y\leftrightarrow v$ (exchanging the two pairs is not required). Let $\Sym$ be the space of symmetric polynomials $q(x,u)$ of degree at most $s-2$ in each variable. It has the basis $x^\alpha u^\beta+x^\beta u^\alpha$ ($\alpha<\beta$), $x^\alpha u^\alpha$, with $0\le\alpha\le\beta\le s-2$, so
\[
d:=\dim\Sym=\binom s2,\qquad \SSS=\Sym\otimes\Sym,\qquad \dim\SSS=d^2
\]
\cite[equation \texttt{eq:dim}]{OAIKn}. Let $\psi:\Sym\to\R$ be evaluation at $(t_a,t_b)$. It is nonzero, as $\psi(1)=1$. Define
\begin{align*}
\SSS_a&=\{P\in\SSS:P(t_a,t_b,y,v)\equiv0\}=\ker(\psi\otimes\mathrm{id}),\\
\SSS_b&=\{Q\in\SSS:Q(x,u,t_a,t_b)\equiv0\}=\ker(\mathrm{id}\otimes\psi).
\end{align*}
By the within-pair symmetries, every $P\in\SSS_a$ also satisfies $P(t_b,t_a,y,v)\equiv0$, and every $Q\in\SSS_b$ satisfies $Q(x,u,t_b,t_a)\equiv0$.

\begin{lemma}[Dimensions]\label{lem:dim}
$\SSS_a=\ker\psi\otimes\Sym$ and $\SSS_b=\Sym\otimes\ker\psi$. Hence $\dim\SSS_a=\dim\SSS_b=d^2-d=\binom s2^2-\binom s2$. Moreover $\SSS_a\cap\SSS_b=\ker\psi\otimes\ker\psi$ has dimension $(d-1)^2$.
\end{lemma}

\begin{proof}
Choose a basis $f_0,\dots,f_{d-1}$ of $\Sym$ with $\psi(f_0)=1$ and $f_1,\dots,f_{d-1}\in\ker\psi$. Writing $P=\sum c_{jk}f_j\otimes f_k$, we get $(\psi\otimes\mathrm{id})P=\sum_kc_{0k}f_k$, which vanishes exactly when the row $j=0$ is zero. The same applies to the column $k=0$ for $\SSS_b$, and to both together for the intersection.
\end{proof}

The evaluation map
\[
E:\SSS\to\R^{\TT},\qquad E(P)_{(\{i,p\},\{k,q\})}=P(t_i,t_p,t_k,t_q),
\]
is well defined by the within-pair symmetries.

\begin{lemma}[Mixed orthogonality]\label{lem:mixed}
Let $\Delta=(x-u)(y-v)(x-y)(x-v)(u-y)(u-v)$. For $\alpha=(\{i,p\},\{k,q\})\in\TT$ put
\[
a_\alpha=J([ip],[kq])\,w_iw_pw_kw_q\,\Delta(t_i,t_p,t_k,t_q),
\]
which does not depend on the orientations $[ip],[kq]$ chosen. Every $a_\alpha$ is nonzero, and
\begin{equation}\label{eq:mixed}
\sum_{\alpha\in\TT}a_\alpha\,E(P)_\alpha\,E(Q)_\alpha=0\qquad(P\in\SSS_a,\ Q\in\SSS_b).
\end{equation}
\end{lemma}

\begin{proof}
Reversing $[ip]$ negates $J$ and exchanges $x,u$ in $\Delta$. That exchange negates $(x-u)$, swaps $(x-y)$ with $(u-y)$ and swaps $(x-v)$ with $(u-v)$, so it negates $\Delta$. The same holds for $[kq]$. Hence $a_\alpha$ is well defined. It is nonzero because $J\ne0$ on $\TT$, all $w_i\ne0$, and the four labels are distinct.

Let $H=PQ\Delta$. The factor $\Delta$ has degree $3$ in each variable, so $H$ has degree at most $2(s-2)+3=2s-1=n-2$ in each variable. Both sides of the first identity in~\eqref{eq:testconditions} vanish because $P\in\SSS_a$, and both sides of the second vanish because $Q\in\SSS_b$. Lemma~\ref{lem:tests} gives $\Sigma(H)=0$. We sort the terms of~\eqref{eq:testsum}. Terms with $i=p$ or $k=q$ vanish since $[ii]=0$. Terms in which $\{i,p\}$ or $\{k,q\}$ is the missing pair vanish by the zero extension. Terms in which the two pairs share a vertex vanish since $\Delta$ does. The remaining terms have four distinct indices and correspond to ordered pairs $\alpha$ of independent edges of $G$, each with four choices of orientations. For each choice the term equals $a_\alpha E(P)_\alpha E(Q)_\alpha$: the orientation signs of $J$ and $\Delta$ cancel as above, and $P,Q$ are symmetric within pairs. Pairs outside $\TT$ have $J=0$. Hence $0=\Sigma(H)=4\sum_{\alpha\in\TT}a_\alpha E(P)_\alpha E(Q)_\alpha$.
\end{proof}

The form $\langle X,Y\rangle=\sum_{\alpha\in\TT}a_\alpha X_\alpha Y_\alpha$ on $\R^\TT$ is symmetric and diagonal with nonzero entries, hence nondegenerate. It is not the skew form $J$. Exchanging the two edges of $\alpha$ gives a different coordinate whose coefficient is $-a_\alpha$, because $J$ is skew while $\Delta(y,v,x,u)=\Delta(x,u,y,v)$. This does not affect nondegeneracy. Lemma~\ref{lem:mixed} says that $E(\SSS_b)\subseteq E(\SSS_a)^\perp$. Neither image needs to be isotropic.

\subsection{The twisted detector}

\begin{lemma}[Twisted detector]\label{lem:twisted}
Let $s\ge2$, let $u_1,\dots,u_{2s}$ be distinct reals with any prescribed linear order, let $M$ be the pairing~\eqref{eq:M}, and let $c\in\{1,\dots,2s\}$. If $A\in\PP_{s-2,s-2}$ satisfies $M\bigl(A,(v-u_c)B'\bigr)=0$ for every $B'\in\PP_{s-1,s-2}$, then $A=0$.
\end{lemma}

\begin{proof}
The value $M(A,B)$ depends only on the product $AB$. Hence $M(A,(v-u_c)B')=M(B',\tilde A)$ with $\tilde A=(v-u_c)A\in\PP_{s-2,s-1}$. By Premise~\ref{U:detector}, the pairing on $\PP_{s-1,s-2}\times\PP_{s-2,s-1}$ is nonsingular. Since $M(B',\tilde A)=0$ for every $B'\in\PP_{s-1,s-2}$, we get $\tilde A=0$. As $\R[u,v]$ is an integral domain and $v-u_c\ne0$, $A=0$.
\end{proof}

The proof does not claim that multiplication by $v-u_c$ is invertible on node values; it is not, since it vanishes at the nodes with $v=u_c$. The lemma uses the nonsingularity of the full rectangular pairing. The square consequence \cite[equation \texttt{eq:detector}]{OAIKn}, which tests $A\in\PP_{s-2,s-2}$ against all of $\PP_{s-1,s-1}$, does not apply directly, because the twisted test functions $(v-u_c)B'$ form a proper subspace of $\PP_{s-1,s-1}$.

\subsection{Injectivity}

\begin{lemma}[Injectivity]\label{lem:injective}
The restrictions of $E$ to $\SSS_a$ and to $\SSS_b$ are injective.
\end{lemma}

\begin{proof}
Let $P\in\SSS_a$ with $E(P)=0$. The proof has three steps.

\emph{Step 1: localization at complete stars.} Let $U(x)$ have degree at most $s$, $V'(y)$ degree at most $s-1$, and $B'\in\PP_{s-1,s-2}$ (in the variables $u,v$). Consider
\begin{equation}\label{eq:loctest}
H=P(x,u,y,v)\,(u-v)(u-y)(v-x)\,U(x)\,(y-t_a)V'(y)\,(v-t_a)B'(u,v).
\end{equation}
Its separate degrees are at most
\[
\begin{array}{lll}
x:&(s-2)+1+s&=2s-1,\\
u:&(s-2)+2+(s-1)&=2s-1,\\
y:&(s-2)+1+1+(s-1)&=2s-1,\\
v:&(s-2)+2+1+(s-2)&=2s-1,
\end{array}
\]
where the middle terms count, in order, the factors $(v-x)$; $(u-v),(u-y)$; $(u-y),(y-t_a)$; and $(u-v),(v-x),(v-t_a)$. All are at most $n-2$. Both sides of the first identity in~\eqref{eq:testconditions} vanish because $P\in\SSS_a$. For the second, $H(x,u,t_a,t_b)=0$ through the factor $(y-t_a)$ and $H(x,u,t_b,t_a)=0$ through the factor $(v-t_a)$. Lemma~\ref{lem:tests} gives $\Sigma(H)=0$.

We evaluate $\Sigma(H)$. Terms with two independent edges of $G$ contribute $J\cdot P(\cdots)\cdots$, which vanishes either because $J=0$ or because the pair lies in $\TT$ and $E(P)$ vanishes there. Terms involving the missing pair or a loop vanish. The factor $(u-v)(u-y)(v-x)$ kills the coincidences $p=q$, $p=k$ and $q=i$. The only surviving terms have $i=k$ with $i,p,q$ distinct, and there
\[
(u-v)(u-y)(v-x)\big|_{(t_i,t_p,t_i,t_q)}=(t_p-t_q)(t_p-t_i)(t_q-t_i).
\]
With $F_i(t)=\prod_{j\ne i}(t-t_j)$ we have $F'(t_p)=(t_p-t_i)F_i'(t_p)$ for $p\ne i$, so $w_p(t_p-t_i)=1/F_i'(t_p)$. Therefore
\begin{equation}\label{eq:localized}
0=\Sigma(H)=\sum_{i=1}^nw_i^2\,U(t_i)(t_i-t_a)V'(t_i)\,\gamma_i(B'),
\end{equation}
where
\[
\gamma_i(B')=\sum_{\substack{p,q\ne i}}J([ip],[iq])\,\frac{t_p-t_q}{F_i'(t_p)F_i'(t_q)}\,P(t_i,t_p,t_i,t_q)\,(t_q-t_a)\,B'(t_p,t_q).
\]

Let $i\notin\{a,b\}$. Every pair $\{i,p\}$ with $p\ne i$ is then an edge of $G$. For $p\ne q$ the edges $ip,iq$ share only the endpoint $i$, so they do not cross (the drawing is normalized), and both orientations $[ip],[iq]$ point out of $i$. By~\eqref{eq:J}, $J([ip],[iq])=\frac12T_i(ip,iq)=:S_i(p,q)$, the sign function of a linear order on the $2s$ indices $p\ne i$; the diagonal terms vanish because of the factor $t_p-t_q$. Hence $\gamma_i(B')=M_i\bigl(P(t_i,u,t_i,v),(v-t_a)B'(u,v)\bigr)$, where $M_i$ is the pairing~\eqref{eq:M} on the $2s$ nodes $\{t_p:p\ne i\}$ with this order and $\Gamma=F_i$. For $i\in\{a,b\}$ the star of $i$ in $G$ is incomplete; we do not interpret $\gamma_a,\gamma_b$ at all.

\emph{Isolation.} Fix $B'$ and $i_0\notin\{a,b\}$. Identity~\eqref{eq:localized} holds for every $U$ and $V'$, and it is linear in their product. The products $U(t)V'(t)$ with $\deg U\le s$ and $\deg V'\le s-1$ span all polynomials of degree at most $2s-1$, because every exponent $k\le 2s-1$ splits as $k_1+k_2$ with $k_1\le s$, $k_2\le s-1$. So we may replace $U(t)(t-t_a)V'(t)$ by any polynomial of degree at most $2s$ that vanishes at $t_a$. The Lagrange polynomial $L_{i_0}$, of degree $n-1=2s$, equal to $1$ at $t_{i_0}$ and $0$ at all other nodes, is such a polynomial because $a\ne i_0$. It also vanishes at $t_b$. Substituting it into~\eqref{eq:localized} kills every term except $i=i_0$, including the uninterpreted terms $i=a,b$, and leaves $w_{i_0}^2\gamma_{i_0}(B')=0$. As $w_{i_0}\ne0$,
\[
M_{i_0}\bigl(P(t_{i_0},u,t_{i_0},v),(v-t_a)B'\bigr)=0\qquad\text{for every }B'\in\PP_{s-1,s-2}.
\]
Here $P(t_{i_0},u,t_{i_0},v)\in\PP_{s-2,s-2}$, and $t_a$ is one of the $2s$ nodes $\{t_p:p\ne i_0\}$. Lemma~\ref{lem:twisted} with $u_c=t_a$ gives $P(t_{i_0},u,t_{i_0},v)=0$ as a polynomial, for each of the $2s-1$ indices $i_0\notin\{a,b\}$.

\emph{Step 2: collision factors.} The polynomial $P(x,u,x,v)$ has degree at most $2s-4$ in $x$. Each of its coefficients in $(u,v)$ vanishes at the $2s-1>2s-4$ distinct values $t_{i_0}$, $i_0\notin\{a,b\}$, so $P(x,u,x,v)=0$. Viewing $P$ as a polynomial in $y$ over $\R[x,u,v]$, we get $(x-y)\mid P$. Applying the symmetries $x\leftrightarrow u$ and $y\leftrightarrow v$ gives divisibility by $u-y$, $x-v$ and $u-v$ as well. Each of these four linear forms is prime in $\R[x,u,y,v]$, since the quotient by it is a polynomial ring, and they are pairwise nonassociate. Since the ring is a unique factorization domain, their product
\[
C=(x-y)(x-v)(u-y)(u-v)
\]
divides $P$.

\emph{Step 3: descent.} We have shown that $C$ divides every element of $\ker E\cap\SSS_a$. Suppose this kernel contains a nonzero element and choose one, $P$, of minimum degree in $x$. Write $P=CR$ with $R\ne0$. Since $C$ has degree exactly $2$ in each variable, $R$ has degree $\deg_zP-2$ in each variable $z$. In particular every separate degree of $P$ is at least $2$. This is already impossible when $s=3$, where all separate degrees are at most $1$. For $s\ge4$, $R$ has separate degrees at most $s-4\le s-2$. The polynomial $C$ is invariant under both within-pair swaps; applying a swap to $P=CR$ and cancelling $C$ in the integral domain shows that $R$ is symmetric within pairs. So $R\in\SSS$. Next,
\[
0=P(t_a,t_b,y,v)=C(t_a,t_b,y,v)\,R(t_a,t_b,y,v),
\]
where $C(t_a,t_b,y,v)=(t_a-y)(t_a-v)(t_b-y)(t_b-v)\ne0$ in $\R[y,v]$. This is a nonzero polynomial, although it vanishes at some points, so $R(t_a,t_b,y,v)\equiv0$ and $R\in\SSS_a$. Finally $C$ is nonzero at every evaluated pair of independent edges, since the four labels are distinct, so $E(R)=0$. Thus $R$ is a nonzero element of $\ker E\cap\SSS_a$ of smaller $x$-degree, a contradiction. Hence $E|_{\SSS_a}$ is injective.

\emph{The second slot.} For $Q\in\SSS_b$ put $Q^{\mathrm{sw}}(x,u,y,v)=Q(y,v,x,u)$. Then $Q^{\mathrm{sw}}\in\SSS_a$, and $E(Q)_{(f,g)}=E(Q^{\mathrm{sw}})_{(g,f)}$ for all $(f,g)\in\TT$. Since $\TT$ is closed under exchanging its two edges, $E(Q)=0$ implies $E(Q^{\mathrm{sw}})=0$, hence $Q^{\mathrm{sw}}=0$ and $Q=0$.
\end{proof}

\subsection{The bound}

\begin{proposition}[Odd lower bound]\label{prop:odd}
For every $s\ge2$, every ordinary drawing of $K_{2s+1}-e$ has at least $\binom s2^2-\binom s2$ crossings.
\end{proposition}

\begin{proof}
For $s=2$ the bound is $0$. Let $s\ge3$ and use the setting of this section. Put $W_a=E(\SSS_a)$ and $W_b=E(\SSS_b)$. By Lemma~\ref{lem:mixed}, $W_b\subseteq W_a^\perp$ for the nondegenerate form $\langle\cdot,\cdot\rangle$ on $V=\R^\TT$. Nondegeneracy identifies $V$ with $V^*$. Restriction of functionals from $V$ to $W_a$ is onto $W_a^*$, and its kernel corresponds to $W_a^\perp$, so $\dim W_a^\perp=|\TT|-\dim W_a$. This holds whether or not the form is nondegenerate on $W_a$. Therefore $\dim W_a+\dim W_b\le|\TT|$. By Lemmas~\ref{lem:injective} and~\ref{lem:dim}, the left side is $2\bigl(\binom s2^2-\binom s2\bigr)$. With~\eqref{eq:support},
\[
\crn(K_{2s+1}-e)=c(D)\ge\frac{|\TT|}2\ge\binom s2^2-\binom s2 .
\]
Every ordinary drawing has at least $c(D)$ crossings.
\end{proof}

\begin{remark}[Why two spaces]\label{rem:split}
For $P,Q\in\SSS$ the polynomial $PQ\Delta$ generally violates~\eqref{eq:testconditions}, so the full space of~\cite{OAIKn} cannot be used. A single admissible space for both slots would be $\SSS_a\cap\SSS_b$. By Lemma~\ref{lem:dim} it has dimension $(d-1)^2$, and its self-orthogonality would give only $\crn\ge(d-1)^2$, short of the target $d^2-d$ by $d-1=\binom s2-1$. Pairing two different spaces of codimension $d$ gives exactly the conjectured value.
\end{remark}

\section{The even lower bound}\label{sec:even}

The following is~\cite[Theorem~3.2]{OWH} specialized to even order with Premise~\ref{U:HH}. Equivalent statements are~\cite[Theorem~3.3]{OWH} and~\cite[Theorem~2]{ChiaLee}. We include the proof so that the paper is complete and to record that the local counts $m_i$ below cancel exactly, so no bound on them is needed. This observation stays within the argument of~\cite{OWH}.

\begin{proposition}\label{prop:even}
For every $M\ge3$, every ordinary drawing of $K_{2M}-e$ has at least $\frac14(M^2-1)(M-2)^2$ crossings.
\end{proposition}

\begin{proof}
Let $n=2M\ge6$, let $x,y$ be the ends of the missing edge and $v_1,\dots,v_{n-2}$ the other vertices, and let $D$ be a normalized minimum drawing of $K_n-xy$ (Premise~\ref{U:norm}). Let $K^x$ and $K^y$ be the subdrawings obtained by deleting $y$ and $x$, respectively. Both are drawings of $K_{n-1}$, so by Premise~\ref{U:HH}
\[
c(K^x)=Z(n-1)+l_1,\qquad c(K^y)=Z(n-1)+l_2,\qquad l_1,l_2\ge0.
\]

\emph{Crossing count by type.} Call the edges among the $v_i$ the core, and the edges at $x$ and at $y$ the two stars. For each $i$ let $a_i$ be the number of crossings on the two edges $xv_i$ and $v_iy$; they do not cross each other. Every crossing of $D$ is between independent edges. A core--core crossing is counted twice in $c(K^x)+c(K^y)$ and not in $\sum a_i$. A crossing between a star edge and a core edge is counted once in $c(K^x)+c(K^y)$ and once in $\sum a_i$. A crossing between $xv_i$ and $yv_j$ ($i\ne j$) is counted in neither subdrawing and twice in $\sum a_i$, once in $a_i$ and once in $a_j$. Two edges of the same star share an endpoint and do not cross. Hence
\begin{equation}\label{eq:typecount}
2c(D)=c(K^x)+c(K^y)+\sum_{i=1}^{n-2}a_i .
\end{equation}

\emph{Restoring the edge.} For each $i$, draw a new arc from $x$ to $y$ that runs alongside $xv_i$, passes around $v_i$ inside a small disk about $v_i$, and then runs alongside $v_iy$. A sufficiently close parallel track follows the polygonal edges through their crossing neighbourhoods. It meets every other edge once for each crossing of that edge with $xv_i$ or $v_iy$, and it does not meet $xv_i$, $v_iy$ or itself. Inside the disk it crosses some number $m_i\ge0$ of the other spokes at $v_i$. Generic local choices avoid vertices, tangencies and triple points. The result is an ordinary drawing $D_i$ of $K_n$ with
\[
c(D_i)=c(D)+a_i+m_i=Z(n)+c_i,\qquad c_i\ge0
\]
by Premise~\ref{U:HH}. The drawing $D_i$ may have crossings between adjacent edges (for instance between the new edge and an edge at $y$ that crossed $xv_i$); this does not affect the lower bound $c(D_i)\ge Z(n)$. Summing over $i$ and eliminating $\sum a_i$ with~\eqref{eq:typecount} gives
\begin{equation}\label{eq:even1}
n\,c(D)=2Z(n-1)+l_1+l_2+(n-2)Z(n)+\sum_ic_i-\sum_im_i .
\end{equation}

\emph{Responsibilities.} Let $r_D(v)$ be the number of crossings on edges incident to $v$. In $D$ each crossing has four distinct endpoints, so $\sum_vr_D(v)=4c(D)$. The crossings of $D$ that involve no edge at $x$ are exactly those of $K^y$, and each other crossing involves exactly one edge at $x$; so $r_D(x)=c(D)-c(K^y)$, and similarly $r_D(y)=c(D)-c(K^x)$. In $D_i$ the old spokes at $v_i$ still do not cross each other, the new edge is not incident to $v_i$, and the only new crossings on spokes at $v_i$ are the $m_i$ local ones. Edges crossing $xv_i$ or $v_iy$ in $D$ are not incident to $v_i$. Hence
\[
r_D(v_i)+m_i=r_{D_i}(v_i)=c(D_i)-c(D_i-v_i)\le Z(n)+c_i-Z(n-1),
\]
because $D_i-v_i$ is a drawing of $K_{n-1}$. Summing,
\[
4c(D)\le 2c(D)-2Z(n-1)-l_1-l_2+(n-2)\bigl(Z(n)-Z(n-1)\bigr)+\sum_ic_i-\sum_im_i,
\]
that is,
\begin{equation}\label{eq:even2}
2c(D)\le(n-2)Z(n)-nZ(n-1)-(l_1+l_2)+\sum_ic_i-\sum_im_i .
\end{equation}

\emph{Elimination.} By~\eqref{eq:even1}, $\sum c_i-\sum m_i=n\,c(D)-(n-2)Z(n)-2Z(n-1)-l_1-l_2$. Substituting into~\eqref{eq:even2} gives
\[
(n-2)\,c(D)\ge(n+2)Z(n-1)+2(l_1+l_2)\ge(n+2)Z(n-1).
\]
The $m_i$ have cancelled. With $Z(2M-1)=\frac14(M-1)^2(M-2)^2$,
\[
c(D)\ge\frac{2M+2}{2M-2}\cdot\frac{(M-1)^2(M-2)^2}{4}=\frac{(M^2-1)(M-2)^2}{4}. \qedhere
\]
\end{proof}

\section{The upper bound}\label{sec:upper}

The upper bound is due to Zheng et al.~\cite{ZLYC} (as quoted in~\cite[Lemma~2.3]{OWH}), and Chia and Lee give a cylindrical construction~\cite[Lemma~1]{ChiaLee}. We verify it in the drawing of Premise~\ref{U:upper}.

\begin{proposition}\label{prop:upper}
For every $n\ge4$ there is an ordinary drawing of $K_n-e$ with exactly $Z(n)-\binom{\lfloor(n-1)/2\rfloor}{2}$ crossings.
\end{proposition}

\begin{proof}
Take the drawing of Premise~\ref{U:upper} and delete the edge $\{0,t\}$, $t=\lfloor(n-1)/2\rfloor$. The result is an ordinary drawing of $K_n-e$, whose crossings are those of $K_n$ not involving $\{0,t\}$. An edge $\{i,j\}$ crosses $\{0,t\}$ exactly when it lies on the same side and its endpoints alternate with $0<t$, that is, $0<i<t<j\le n-1$.

If $n=2s+1$, then $m=s$ and $t=s$. The residue of $0+s$ is $s\ge m$, so $\{0,s\}$ lies below, where the residues are $s,\dots,2s$. For $1\le i\le s-1$ and $s+1\le j\le 2s$ we have $s+2\le i+j\le 3s-1$. If $i+j\le2s$ the residue is $i+j\in[s+2,2s]$ (below); otherwise it is $i+j-(2s+1)\in[0,s-2]$ (above). The crossing edges are those with $j\le 2s-i$, and their number is $\sum_{i=1}^{s-1}(s-i)=\binom s2$.

If $n=2s$, then $m=s$ and $t=s-1$. The residue $s-1<m$ puts $\{0,s-1\}$ above, where the residues are $0,\dots,s-1$. For $1\le i\le s-2$ and $s\le j\le2s-1$ we have $s+1\le i+j\le 3s-3$. If $i+j\le 2s-1$ the residue is $i+j\ge s+1$ (below); otherwise it is $i+j-2s\in[0,s-3]$ (above). The crossing edges are those with $j\ge2s-i$, and their number is $\sum_{i=1}^{s-2}i=\binom{s-1}2$.

In both cases the number of crossings on $\{0,t\}$ is $\binom t2$, and the remaining drawing has $Z(n)-\binom t2$ crossings.
\end{proof}

\section{Proofs of the main results}\label{sec:proofs}

\begin{proof}[Proof of Theorem~\ref{thm:A}]
By edge-transitivity we may fix the missing edge. For $n=2,3$ the graph $K_n-e$ is a forest and the formula gives $0$. For $n\ge4$ Proposition~\ref{prop:upper} gives the upper bound. For the lower bound, odd $n=2s+1\ge5$ is Proposition~\ref{prop:odd}, $n=4$ has target $0$, and even $n=2M\ge6$ is Proposition~\ref{prop:even}. The parity forms~\eqref{eq:parity} and the floor form were matched in Section~\ref{sec:intro}.
\end{proof}

\begin{proof}[Proof of Corollary~\ref{cor:B}]
By~\cite[Corollary~5.1]{OWH}, for $m\ge4$ and $n\ge1$, $\crn(K_m\,\square\,P_n)=(n-1)\crn(K_{m+2}-e)+2\crn(K_{m+1})$. We insert Theorem~\ref{thm:A} and Premise~\ref{U:HH}.

If $m=2k$ ($k\ge2$), then $\crn(K_{2k+2}-e)=\frac14k(k+2)(k-1)^2$ and $\crn(K_{2k+1})=\frac14k^2(k-1)^2$, so
\[
\crn(K_m\,\square\,P_n)=\frac{k(k-1)^2}{4}\bigl((n-1)(k+2)+2k\bigr)=\frac{k(k-1)^2}{4}\bigl(n(k+2)+k-2\bigr).
\]
The floors in~\eqref{eq:zheng} are $k,\,k-1,\,k-1,\,k+2,\,k-2$.

If $m=2k+1$ ($k\ge2$), then $\crn(K_{2k+3}-e)=\frac14(k+2)(k+1)k(k-1)$ and $\crn(K_{2k+2})=\frac14(k+1)k^2(k-1)$, so
\[
\crn(K_m\,\square\,P_n)=\frac{(k+1)k(k-1)}{4}\bigl((n-1)(k+2)+2k\bigr)=\frac{(k+1)k(k-1)}{4}\bigl(n(k+2)+k-2\bigr).
\]
The floors in~\eqref{eq:zheng} are $k+1,\,k,\,k-1,\,k+2,\,k-2$. For $n=1$ the minus-edge term drops out and the value is $2Z(m+1)$; for $m=4$ the value is $2n$.
\end{proof}

\section{Novelty boundary}\label{sec:novelty}

\paragraph{Routine ranges.} For odd $n=2s+1$, inserting Premise~\ref{U:HH} into~\cite[Theorem~3.1]{OWH},
\[
\crn(K_n-e)\ge\frac1n\Bigl(2Z(n-1)+(n-2)\bigl(Z(n)-\lfloor\tfrac{n-3}2\rfloor\bigr)\Bigr),
\]
gives a bound exactly $\frac{(s-1)(s-2)}{2(2s+1)}$ below the target. Since $\frac{(s-1)(s-2)}{2(2s+1)}-1=\frac{s(s-7)}{2(2s+1)}$, the deficit is less than $1$ exactly for $s\le6$. So integrality closes $n\le13$ (for example the bound is $2720/13$ at $n=13$, with ceiling $210$) and fails for every odd $n\ge15$ (at $n=15$ the ceiling is $419$, the target $420$). Other classical routes also stay below the target for odd $n\ge15$. Evaluating~\cite[Theorem~3.2]{OWH} at odd order leaves a gap of $\frac{s(s-1)(s-2)}{2(2s-1)}$. Averaging over vertex deletions, charging $Z(2s)$ for the two ends of the missing edge and the even value of Theorem~\ref{thm:A} for the others, leaves a gap of $\frac{(s-1)^2(s-2)}{2(2s-3)}$. These identities are checked exactly in Section~\ref{sec:computations}. They delimit what these particular routes give; they are not a proof that no simpler argument exists.

\paragraph{Status before this work.} The surveys~\cite{CHN} (Section~2.2.1 and Section~3.2) list the minus-edge conjecture as known for $n\le12$ and the product conjecture as confirmed for $m\le10$. The OpenAI preprint~\cite{OAIKn} proves only $\crn(K_n)$; its companion~\cite{OAIKmn} treats $\crn(K_{m,n})$. Neither states a result on $K_n-e$ or on Cartesian products; we checked the statements, the final sections, and the README files at the cited commit.

\paragraph{Comparison with Sadhu's bipartite theorem.} Sadhu~\cite[Theorem~1.1]{Sadhu} proves an exact formula, written there as $\crn(K_{m,n}-e)=Z(m,n)-t_mt_n$, for all positive $m,n$ and every edge $e$, and states that the argument does not assume the Zarankiewicz formula. We read the preprint in the text rendering of its ResearchGate page, because the PDF link was not accessible to us. The passages compared are the abstract, the contribution paragraph, the balanced compression inequality (Lemma~3.2) and the beginning of its proof, the rectangular bipartite detector (Lemma~4.1), the signed three-edge completion of the missing edge in the proof of Theorem~1.1, the closing rank comparison, and the closing questions. In those passages the missing edge is replaced by a signed combination of three present edges and the intersection identity is run on a completed matrix. A multi-index compression inequality absorbs the dimension losses. Our argument instead forces the test chains to have zero coefficient on the missing pair, pairs two different vanishing subspaces of one four-variable space, quotes the already proved rectangular nonsingularity for complete graphs, and never localizes at the incomplete stars. We have not re-proved Sadhu's lemmas, and our comparison is limited to the passages listed.

\paragraph{Limits.} Our literature search was bounded: the surveys above, the primary papers~\cite{OWH,ChiaLee,Asiri}, and targeted web searches in October 2026. We found no earlier proof of the odd case $n\ge15$, but this is not a priority certificate.

\section{Computational support}\label{sec:computations}

The proofs above do not depend on any computation. The accompanying package contains finite checks that test identities, bookkeeping and hypotheses used in the proof. They use several kinds of arithmetic, listed in Table~\ref{tab:checks}. Not all of them are exact rational computations.

\begin{table}[ht]
\centering
\small
\begin{tabular}{@{}P{0.17\textwidth} P{0.45\textwidth} P{0.32\textwidth}@{}}
\toprule
Check & Content & Arithmetic \\
\midrule
R1--R5 & Closed forms ($n\le3000$), two-page deletion count ($n\le56$), even identity, routine-range gaps (odd $n\le401$), Corollary~\ref{cor:B} algebra ($m\le400$, $n\le12$) & Exact integers (integer division only) and rationals; the two-page count is combinatorial \\
R6 & Rectangular detector instances and twisted detector of Lemma~\ref{lem:twisted}, $s\le5$, random orders, every twist node & Exact rational ranks \\
R7 & $\dim\SSS_a=\dim\SSS_b=d^2-d$, $\dim(\SSS_a\cap\SSS_b)=(d-1)^2$, $s\le6$ & Exact rational ranks \\
R8 & Degree budgets and the identities~\eqref{eq:testconditions} for the tests of Lemmas~\ref{lem:mixed} and~\ref{lem:injective}, $3\le s\le7$; enforced controls showing each twist factor is needed and the untwisted test of~\cite{OAIKn} is inadmissible & Exact polynomial identities over $\mathbb Q$ \\
R9 & Four-orientation factor and the non-converse $H=xy$, for random skew forms (not drawings), $n=7,9$ & Exact rationals \\
Author checker & Lemma~\ref{lem:tests}, orthogonality and ranks on two-page and random straight-line drawings, $n\le12$ & Rationals for identities; floating-point coordinates and angles for straight-line geometry; ranks modulo $2^{61}-1$ \\
Author mod-$p$ checker & Ranks and orthogonality for $13\le n\le19$ (retained log); $n=15,17,19$ re-executed for this package & NumPy arithmetic modulo $2^{31}-1$ \\
Reviewer checker & Independent controls, including $n=15$ & Exact integer geometry, rationals, ranks modulo $1000003$ \\
\bottomrule
\end{tabular}
\caption{Finite checks. R1--R9 are new for this package; the other three are unchanged copies of earlier internal scripts.}
\label{tab:checks}
\end{table}

The release runner fails if any check fails. This matters for the author's mod-$p$ script, which prints a failure message but still exits with status zero. The runner treats that message as a failure, and a self-test confirms it. The author's stdlib checker prints the values of tests that violate~\eqref{eq:testconditions} but does not assert anything about them; the enforced negative controls are R8 and R9. Early runs of the independent reviewer's checker failed because the harness wrongly required certain invalid tests to be nonzero. For example, at $n=15$ the first chosen missing edge had no crossings in the full straight-line drawing. These were faulty expectations in the harness, not counterexamples. The final checker distinguishes them, and its retained logs are cited in the package rather than reproduced here. No finite check establishes the theorem for all $n$.

\section{Formal verification status}\label{sec:formal}

No machine-checked proof of Theorem~\ref{thm:A}, Corollary~\ref{cor:B}, or any lemma of this paper exists at the time of writing, and none is claimed. A separate formalization effort is in progress. Its intended scope, which has not been implemented or audited, is the algebraic core: the tensor-kernel step of Lemma~\ref{lem:tests}, Lemma~\ref{lem:dim}, the twisted detector of Lemma~\ref{lem:twisted} from an explicit statement of Premise~\ref{U:detector}, the degree and missing-pair identities, the root count, division and descent of Lemma~\ref{lem:injective}, the dimension bound in Proposition~\ref{prop:odd}, and the arithmetic of Sections~\ref{sec:even}--\ref{sec:proofs}. The drawing infrastructure (Premises~\ref{U:norm}, \ref{U:cycles}, \ref{U:HH} and~\ref{U:upper}) would enter as explicit hypotheses. This paragraph will be revised only after an actual formalization has been completed and independently audited.

\section*{Acknowledgements and use of AI assistance}
This work was carried out with substantial assistance from AI systems, and the author takes responsibility for its content. The proof of the odd case was first drafted by an AI research agent (Claude Opus~5.5, run in the Factory Droid harness). A separate AI agent (GPT-6.1-Sol) carried out an adversarial proof review and found no required mathematical repair. A third AI agent (Grok~4.7) carried out a contribution and literature review. These internal reviews are not peer review. The manuscript was typeset with the help of an AI agent (Claude Opus~5.5). The mathematics of~\cite{OAIKn} is due to its authors, and the results of~\cite{OWH,ChiaLee,ZLYC,Sadhu} are due to theirs.

\paragraph{License.} The text of this paper is available under the Creative Commons Attribution 4.0 International license. Accompanying original code is under the MIT license. Third-party works keep their own licenses.

\begingroup
\small
\raggedright
\begin{thebibliography}{99}

\bibitem{Asiri}
Z.~Asiri, R.~Burdett, M.~Chimani, M.~Haythorpe, A.~Newcombe, and M.~H. Wagner,
\emph{A systematic approach to crossing numbers of Cartesian products with paths},
in: 33rd International Symposium on Graph Drawing and Network Visualization (GD 2025), LIPIcs vol.~357, article~20, 2025.
\doi{10.4230/LIPIcs.GD.2025.20}.

\bibitem{Bokal2007}
D.~Bokal,
\emph{On the crossing numbers of Cartesian products with paths},
J. Combin. Theory Ser. B \textbf{97} (2007), 381--384.
Cited as reference~[5] of~\cite{OWH}; not consulted directly.

\bibitem{ChiaLee}
G.~L. Chia and C.~L. Lee,
\emph{Crossing numbers of nearly complete graphs and nearly complete bipartite graphs},
Ars Combin. \textbf{121} (2015), 437--446.

\bibitem{CHN}
K.~Clancy, M.~Haythorpe, and A.~Newcombe,
\emph{A survey of graphs with known or bounded crossing numbers},
arXiv:1901.05155v5 (version~2.0, 8~December 2021).
\url{https://arxiv.org/abs/1901.05155v5}.

\bibitem{HararyHill1963}
F.~Harary and A.~Hill,
\emph{On the number of crossings in a complete graph},
Proc. Edinburgh Math. Soc. (2) \textbf{13} (1963), 333--338.

\bibitem{OAIKmn}
OpenAI,
\emph{The crossing number of complete bipartite graphs},
OpenAI Math Release preprint, 23 September 2026.
\url{https://github.com/openai/math}, commit \texttt{fd4aeeb2ee4fc729c18d98444fed42fd0529eeeb}, directory \path{preprints/The-crossing-number-of-complete-bipartite-graphs-September-23-2026/}.

\bibitem{OAIKn}
OpenAI,
\emph{The crossing number of complete graphs},
OpenAI Math Release preprint, 23 September 2026.
\url{https://github.com/openai/math}, commit \texttt{fd4aeeb2ee4fc729c18d98444fed42fd0529eeeb}, file \path{preprints/The-crossing-number-of-complete-graphs-September-23-2026/build/main.tex}.
Theorem, lemma and equation numbers in the text refer to this source.

\bibitem{OWH}
Z.~Ouyang, J.~Wang, and Y.~Huang,
\emph{The crossing number of the Cartesian product of paths with complete graphs},
Discrete Math. \textbf{328} (2014), 71--78.
\doi{10.1016/j.disc.2014.03.011}.

\bibitem{Sadhu}
A.~Sadhu,
\emph{The crossing number of $K_{m,n}$ with one edge deleted},
preprint, October 2026, ResearchGate publication 415807828, CC~BY~4.0.
\url{https://www.researchgate.net/publication/415807828}.

\bibitem{ZLYC}
W.~Zheng, X.~Lin, Y.~Yang, and C.~Cui,
\emph{On the crossing number of $K_m\,\square\,P_n$},
Graphs Combin. \textbf{23} (2007), 327--336.
\doi{10.1007/s00373-007-0726-z}.
Not consulted directly; see the text.

\end{thebibliography}
\endgroup

\end{document}
